% REV_P0 2026-11-24：由 cellcharter_study_introduction_module_v1_20261122.tex 派生；V1保持不变。
% 本轮只改写与CellCharter原作者已有DLPFC工作的重叠定位；证据见 REV_P0_20261124/REV_P0_EVIDENCE_TABLE_20261124.md。
% CellCharter研究专题引言模块（非整篇硕士毕业论文统一第一章）
% 可以根据导师确定的毕业论文总结构并入总论或专章绪论。
% 本文只覆盖第三篇空间聚类研究；不替代其他科研工作的引言。
\section{研究动机、目的与主要内容}
\label{sec:cc-intro-module-v1}

\subsection{研究动机}
随着空间转录组技术的发展，研究者能够把基因表达与组织采样位置联合分析，进而划分与分子特征和组织结构相关的空间域。然而，空间域划分的质量并不由算法名称单独决定。使用表达还是空间增强表示、引入多少层邻域信息、采用何种聚类数选择规则，以及用哪一类指标评价结果，都会影响聚类分析所得结论。

现有空间聚类框架已经提出多尺度邻域聚合、空间图表示学习与空间先验建模等机制。CellCharter作为其中一种成熟方法，可以在已有单点表示上叠加多层邻域信息\cite{varrone2024}。其原作者已在同一套12张DLPFC切片上完成基准比较：以151507、151672、151673三张切片调参，在其余9张切片上评价ARI；调参网格包括2、3、4层邻域，最终选用4层邻域、5000个高变基因和5维scVI潜变量\cite{varrone2024,cellcharteranalyses2024}。因此，“邻域层数会影响DLPFC聚类ARI”并非本研究的首次观察。但在实际数据分析中，扩大邻域表示范围与增加输入维数常常同时发生；若只比较单一数据集或少量随机初始化，容易将由聚类数、样本差异或标签粒度引起的变化误判成表示机制的必然优势。

空间转录组聚类还存在评价目标不同的问题。ARI和NMI比较预测分区与参考层标注的一致性，却不直接告诉研究者真实层间交界处是否发生边界错位或漏检。对于具有明确皮层分层的DLPFC和组织区域结构复杂的乳腺癌，这类差异在解释模型输出时尤其需要被单独记录。

\subsection{研究目的}
本研究不以提出一种新的空间聚类神经网络为目标，而是将已经公开的CellCharter表示与GMM聚类流程视为分析对象，通过一套尽量可复核的多条件评价，考察如下问题：在固定样本与聚类配置时，邻域深度和特征块组合的变化如何影响聚类分区；在不依赖参考类别数确定K时，上述观察是否稳定；整体ARI/NMI和严格空间边界指标是否提供一致的判断；结果与成熟基线相比处于什么水平，以及在另一组织患者队列上具有怎样的异质性。

由于这些问题涉及不同层级的实验比较，本研究分别定义Oracle-K与K-blind两类K设置，并保留原始预测、技术种子、参考mask及文件来源SHA。所有分析仅以已执行的实验为事实基础；对于研究完成后追加的诊断性评价，单独标注其为事后分析，不能将其当作前瞻性独立验证设计。

\subsection{研究内容}
为回答上述问题，本研究围绕12张DLPFC切片及6例乳腺癌Visium样本开展分阶段实验。

首先，基于已形成的scVI低维表达表示和CellCharter四层空间邻域聚合缓存，在DLPFC上比较5D至25D的累计特征与多种等维15D特征组合，并在10个技术种子上评价聚类结果。在5000个高变基因、5维潜变量条件下，累计15D、20D、25D分别对应原作者调参网格中的2、3、4层邻域，其中25D即原作者单切片基准的最终配置，可用于对照原作者已发表的结果；5D（不聚合）、10D（1层）和非连续的等维15D组合不在原作者网格中。其次，在同一冻结输入上设定固定K网格和无标签AutoK两类补充实验，分析聚类结果对簇数选择方式的敏感性。

再次，利用原有预测及六邻接阵列坐标，对严格空间边界的正确检出、误报和漏检进行重新审查，并分别提供1-hop、2-hop容差的事后敏感性评价。随后，在同一DLPFC切片上纳入非空间GMM、Leiden以及成熟空间方法BANKSY和BayesSpace，明确不同模型的原生数据处理及K选择差别。最后，在6名乳腺癌患者的独立数据上重新执行同类特征及基线流程，检查跨患者效果方向、低表现样本和方法间负结果。

\subsection{研究工作特点与论文组织}
本研究的主要特点是将多尺度特征、聚类数设置、边界评价、技术随机性与独立组织队列放在同一可追溯证据框架内；同时明确区分成熟方法与本文自行完成的实验分析，避免把已有邻域聚合操作写成新的算法结构。相对于原作者的DLPFC工作，本研究的新增内容限于以下几项：在全部12张切片、10个GMM种子上补充0层与1层累计特征以及非连续等维15D块替换（2--4层的比较与原作者调参网格重叠，仅样本范围和重复结构不同）；预设固定K网格；逐切片的无标签AutoK选K及其事后评分；ARI/NMI之外的严格与容差空间边界评价；双K策略下的成熟基线对照；以及乳腺癌6例患者的外部复核。上述内容均为实证评价，不构成新的聚类方法。具体结果并不以全部评价指标同向改善为前提，边界定位、不同患者的表现以及基线的局限都将作为真实结果完整呈现。

本专题的相关方法背景见后续研究背景部分，实际数据与技术路线见研究方法部分，完整结果和负面观察见实验结果与分析部分。上述内容可作为毕业论文中的一个独立研究专题，最终章节编号与其同其他论文工作的衔接由学位论文整体结构确定。
